\section{Disturbances, sensitivity, and prefilters}

\subsection{Disturbance transfer functions}

\subsubsection{Superposition for multiple inputs}
\[
Y(s)=G_{yr}(s)R(s)+G_{yd_1}(s)D_1(s)+G_{yd_2}(s)D_2(s)
\]
Linear systems allow each reference or disturbance contribution to be analysed separately.

\subsubsection{Output disturbance after the plant input}
\[
Y(s)=\frac{G(s)C(s)}{1+C(s)G(s)}R(s)
+\frac{1}{1+C(s)G(s)}D(s)
\]
For an additive output disturbance, feedback suppresses the disturbance by the sensitivity function.

\subsubsection{Input disturbance before the plant}
\[
Y(s)=\frac{G(s)C(s)}{1+C(s)G(s)}R(s)
+\frac{G(s)}{1+C(s)G(s)}D(s)
\]
The plant filters the disturbance before feedback attenuation.

\subsubsection{Measurement disturbance}
\[
Y(s)=\frac{G(s)C(s)}{1+C(s)G(s)}R(s)
-\frac{G(s)C(s)}{1+C(s)G(s)}N(s)
\]
Measurement noise passes through the complementary sensitivity path.

\subsection{Sensitivity functions}

\subsubsection{Sensitivity function}
\[
S(s)=\frac{1}{1+L(s)}
\]
Sensitivity maps output disturbances and model uncertainty effects in standard unity feedback.

\subsubsection{Complementary sensitivity function}
\[
T(s)=\frac{L(s)}{1+L(s)}
\]
Complementary sensitivity maps reference tracking and measurement-noise transmission.

\subsubsection{Input sensitivity}
\[
G_{ur}(s)=\frac{C(s)}{1+C(s)G(s)}
\]
The control input can become large near bandwidth limits, resonance peaks, or fast reference commands.

\subsubsection{Steady-state disturbance error}
\[
e_{\mathrm{ss},d}=\lim_{s\to0}sE_d(s)
\]
The final value theorem gives constant-disturbance error when the closed-loop poles are stable.

\subsubsection{Disturbance magnitude from Bode plot}
\[
|y_{\mathrm{ss}}|=|G_{yd}(j\omega)|\,|d_{\mathrm{ss}}|
\]
At very low frequency, the Bode magnitude approximates constant disturbance sensitivity.

\subsection{Prefilters}

\subsubsection{Reference prefilter}
\[
Y(s)=\frac{F(s)L(s)}{1+L(s)}R(s)
\]
A prefilter shapes reference tracking without changing the loop transfer function and stability margins.

\subsubsection{First-order prefilter}
\[
F(s)=\frac{1}{\tau_f s+1}
\]
A first-order prefilter reduces fast reference content and can reduce overshoot or input peaks.

\subsubsection{Lead prefilter}
\[
F_{\mathrm{lead}}(s)=\frac{\tau_f s+1}{\alpha_f\tau_f s+1}
\]
The slides use lead-shaped prefilters to improve reference response without changing feedback robustness.
